\section{Описание применения генеративной модели}

В ходе выполнения курсового проекта в качестве вспомогательного инструмента использовалась генеративная модель ChatGPT~\cite{GPT}. 

Модель использовалась, во-первых, для выделения ключевых идей в источниках литературы с последующей проверкой фактов и цитат из этих источников. 

Во-вторых, для написания тестовых сценариев и нахождения граничных случаев, при которых программа могла работать некорректно. В частности, модель помогала формулировать проверки для сценариев авторизации, переходов между фазами игры, обработки WebSocket-соединений и ограничений, связанных с WIP-лимитами. 

В-третьих, для краткой суммаризации уже реализованной функциональности другими участниками команды. Последнее особенно важно, так как разработка программы велась параллельно, генеративная модель позваляла быстро ознакомиться с результатами проделанной работы других участников коллектива и быстрее согласовывать отдельные фрагменты документации.












